% ======================================================================
% CAPÍTULO 1 — INTRODUÇÃO
% ======================================================================
\chapter{INTRODUÇÃO}

\section{Contextualização do tema}

O mercúrio (Hg) é um poluente global de elevada toxicidade que pode ser
liberado para o ambiente por fontes naturais e antropogênicas. Seu
comportamento, transporte e risco toxicológico dependem fortemente da espécie
química em que se encontra: o metilmercúrio, por exemplo, é uma forma
neurotóxica que se acumula na cadeia alimentar, enquanto o mercúrio elementar
apresenta elevada pressão de vapor e grande mobilidade atmosférica. Nesse
contexto, a \emph{especiação} de mercúrio --- isto é, a determinação das
diferentes formas químicas do elemento em uma amostra --- é indispensável para
avaliações de risco e estudos biogeoquímicos confiáveis \cite{leermakers2005,
bloom1989}.

A determinação de espécies de mercúrio exige um procedimento analítico de
preparação de amostras composto por etapas sequenciais bem definidas. No sistema
considerado neste trabalho, a amostra líquida é submetida à \emph{derivatização}
com tetraetilborato de sódio (TEBS, NaBEt$_4$), que volatiliza as espécies de
mercúrio presentes. Em seguida, um gás de arraste (hélio) conduz os vapores até
um Tubo U imerso em um copo com nitrogênio líquido, onde ocorre a
\emph{criofocalização} das espécies a aproximadamente $-196\,\celsius$. Após o
período definido pelo operador, o copo é abaixado, o aquecimento do Tubo U é
iniciado e as espécies são liberadas progressivamente conforme sua temperatura
de volatilização, sendo então atomizadas a $700\,\celsius$ em um segundo forno e
detectadas por espectrometria de absorção atômica.

A qualidade analítica do processo depende diretamente da forma como a
temperatura do Tubo U é controlada ao longo do tempo: idealmente, uma rampa
linear de temperatura na faixa de aproximadamente $-50\,\celsius$ a
$230\,\celsius$ permite correlacionar o instante de liberação de cada espécie
com sua temperatura de volatilização. Qualquer não linearidade na rampa
compromete a separação das espécies e produz picos sobrepostos no detector,
invalidando a análise. O processo físico completo é detalhado no Capítulo~2, e
os requisitos técnicos que o especificam estão documentados nos arquivos de
requisitos e de especificação do sistema.

\section{Problema de pesquisa}

O equipamento de preparação de amostras descrito era originalmente controlado
por um sistema supervisório desenvolvido em LabVIEW, executado em um computador
com comunicação serial a um módulo de aquisição de dados (Arduino Uno). Com o
tempo, esse sistema legado passou a apresentar limitações relevantes, que podem
ser agrupadas em cinco frentes principais:

\begin{enumerate}[label=\roman*)]
  \item \textbf{Controle inadequado da rampa de temperatura do Tubo U.} O
  sistema não garante a linearidade da rampa na faixa de $-50\,\celsius$ a
  $230\,\celsius$, comprometendo a separação e a identificação das espécies de
  mercúrio e resultando em picos sobrepostos, análises inválidas e retrabalho;
  \item \textbf{Parâmetros não persistentes e pouco rastreáveis.} Os ganhos dos
  controladores e os tempos do método ficavam fixos ou sem persistência
  confiável entre execuções, prejudicando a repetibilidade e a auditoria do
  método;
  \item \textbf{Interface propensa a erro humano.} A interface legada não
  dispunha de um sinótico claro do processo, de indicadores de tendência ou de
  uma parada de emergência centralizada de alta prioridade;
  \item \textbf{Acoplamento entre a lógica de aplicação e a execução em tempo
  real.} A ausência de separação entre o supervisório e a aquisição de campo
  dificultava a depuração, a manutenção e a evolução do sistema;
  \item \textbf{Segurança insuficiente.} Não havia um mecanismo de watchdog que
  desligasse os aquecedores automaticamente em caso de perda de comunicação,
  expondo o equipamento ao risco de congelamento permanente do Tubo U ou de
  sobreaquecimento dos fornos.
\end{enumerate}

A partir desse diagnóstico, o problema de pesquisa pode ser enunciado da
seguinte forma: \emph{como projetar e implementar um sistema de automação e
supervisão que garanta o controle preciso e seguro do processo de preparação de
amostras em especiação de mercúrio, superando as limitações do sistema legado e
preservando a repetibilidade analítica do método?}

\section{Justificativa}

A substituição do sistema legado é motivada, em primeiro lugar, pela necessidade
de \emph{repetibilidade analítica comprovável}. Métodos de especiação dependem
da reprodutibilidade da rampa de temperatura e da estabilidade das condições do
processo; a persistência dos parâmetros do método e o registro de tendências
permitem auditoria e rastreabilidade das análises. Em segundo lugar, a
modernização apoia-se na disponibilidade de plataformas de baixo custo e alta
confiabilidade --- como o Raspberry Pi e o Arduino Uno --- capazes de atender
aos requisitos de tempo real e de interface do processo a um custo muito menor
que soluções baseadas em controladores lógicos programáveis industriais ou em
licenças proprietárias de software de supervisão \cite{arduino, raspberrypi}.

Em terceiro lugar, a dimensão de \emph{segurança operacional} é crítica em um
ambiente laboratorial que manipula nitrogênio líquido e aquecedores de alta
temperatura. A ausência de uma parada segura automática representa risco de dano
ao equipamento --- como o congelamento permanente do Tubo U --- e à amostra.
Por fim, do ponto de vista acadêmico, o trabalho integra conceitos de controle
de processos, sistemas embarcados, comunicação industrial, engenharia de
software e interface homem-máquina, constituindo um estudo de caso aplicado de
engenharia de automação.

\section{Objetivos}

\subsection{Objetivo geral}

Projetar e implementar um sistema de automação e supervisão para a preparação de
amostras em especiação de mercúrio, baseado em máquina de estados finita,
controle PID misto, módulo de aquisição de dados em tempo real e interface
homem-máquina Web, com parâmetros persistentes e segurança operacional.

\subsection{Objetivos específicos}

\begin{enumerate}[label=\arabic*)]
  \item Modelar o processo de preparação de amostras (fases $T_0$ a $T_3$) em
  uma máquina de estados finita com matriz de acionamento dos atuadores (SV1 a
  SV5, bomba e fornos);
  \item Implementar o controle PID misto do Forno 1 (Tubo U) --- razão de taxas
  de aquecimento para temperaturas abaixo de $0\,\celsius$ e controle PID em
  malha fechada acima desse limiar --- e o controle do Forno 2 (atomizador) em
  setpoint fixo de $700\,\celsius$;
  \item Definir e implementar o protocolo de comunicação JSON estruturado entre
  o Raspberry Pi e o Arduino, operando em ciclo de controle de 4~Hz, com
  handshake de parametrização e pacotes de escrita e de leitura;
  \item Desenvolver o firmware do módulo de aquisição de dados no Arduino Uno,
  responsável pelo acionamento dos atuadores, pela leitura de termopares via
  SPI e pelo watchdog de segurança;
  \item Desenvolver uma IHM Web (React e TypeScript) com monitoramento em tempo
  real, Diagrama de Tempos, gráficos de tendência e painéis de controle e de
  configuração;
  \item Implementar a persistência atômica dos parâmetros do método em arquivo
  JSON, com backup rotativo e validação de faixas;
  \item Validar o sistema por meio de testes unitários, de integração,
  funcionais e de segurança, seguindo a estratégia de desenvolvimento orientado
  a testes adotada no projeto.
\end{enumerate}

\section{Hipótese}

A hipótese de trabalho é que uma arquitetura mestre-escravo composta por um
Raspberry Pi executando a lógica de supervisão e controle em Python e um Arduino
Uno atuando como DAQ em tempo real, integrados por um protocolo JSON a 4~Hz e
supervisionados por uma IHM Web, é capaz de atender aos requisitos funcionais e
de segurança do processo de preparação de amostras, proporcionando controle
preciso da rampa de temperatura, persistência auditável dos parâmetros e parada
segura automática diante de falhas de comunicação, com custo e flexibilidade
superiores aos do sistema legado.

\section{Delimitação do trabalho}

O escopo deste trabalho limita-se à modernização do software de automação e
supervisão do processo de preparação de amostras (versão V1), conforme definido
na documentação de requisitos do sistema. Não fazem parte do
escopo: a integração automática com o detector Lumex (apenas a indicação do
início de leitura é considerada), a calibração automática dos termopares, o
controle de vazão dos rotâmetros (regulagem manual) e mecanismos de
autenticação multiusuário, previstos para versões futuras. A validação
experimental em bancada com o processo físico, incluindo a calibração da curva
de aquecimento do Tubo U e a verificação da linearidade da rampa, é tratada como
etapa de continuidade, sendo o presente desenvolvimento validado em ambiente de
simulação por meio de testes automatizados.

\section{Visão geral da metodologia}

Do ponto de vista metodológico, o trabalho caracteriza-se como uma pesquisa
aplicada, de natureza experimental e abordagem quantitativa, organizada segundo
as etapas de: (i) levantamento e análise dos requisitos do processo a partir da
documentação técnica existente; (ii) definição da arquitetura do sistema; (iii)
desenvolvimento incremental do firmware, do backend e da IHM Web; e (iv)
validação por testes automatizados. O desenvolvimento seguiu práticas de
especificação dirigida e de desenvolvimento orientado a testes, em que os
contratos de interface e os critérios de aceitação foram definidos antes da
implementação \cite{beck2002}. O detalhamento metodológico é
apresentado no Capítulo~4.

\section{Estrutura da monografia}

Este trabalho está organizado em sete capítulos. O Capítulo~1 (presente)
contextualiza o tema e apresenta o problema, a justificativa e os objetivos. O
Capítulo~2 desenvolve a fundamentação teórica, abordando a química do processo,
os conceitos de controle e a base tecnológica empregada. O Capítulo~3 apresenta
os trabalhos relacionados e o posicionamento da solução proposta. O Capítulo~4
descreve a metodologia adotada. O Capítulo~5 detalha o desenvolvimento do
sistema, desde a arquitetura até a implementação de cada módulo. O Capítulo~6
apresenta os experimentos e os resultados obtidos na validação, com a respectiva
discussão. Por fim, o Capítulo~7 apresenta as conclusões e as sugestões para
trabalhos futuros.
